\chapter{Fazit und Ausblick}

\section{Beantwortung der Leitfrage und der Projektfragen}

Die Leitfrage dieser Arbeit lautet, ob sich ein kostengünstiges autonomes mobiles Robotersystem so auslegen lässt, dass Fahrkern, Sensor- und Sicherheitsbasis, Lokalisierung und Kartierung, Navigation sowie Bedien- und Leitstandsebene im Innenraum reproduzierbar zusammenarbeiten. Die Ergebnisse aus Kapitel 6 beantworten diese Leitfrage im Kern positiv. Der Prototyp erreicht eine funktionsfähige und sicher priorisierte Innenraum-Mobilität auf Basis einer verteilten Architektur aus zwei ESP32-S3-Mikrocontrollern und einem Raspberry Pi 5. Gleichzeitig zeigt die Validierung, dass nicht alle Akzeptanzkriterien in gleicher Tiefe nachgewiesen sind und dass erweiterte Funktionen nicht mit dem gleichen Reifegrad vorliegen wie der fahrkritische Kern.

Die Arbeit zielt damit nicht auf den vollständigen Nachweis industrieller Serienreife, sondern auf den belastbaren Nachweis einer tragfähigen Systemarchitektur für einen kostengünstigen Innenraum-AMR. Diese Einordnung entspricht auch der Roadmap: Zuerst müssen Fahrkern, Sensor- und Sicherheitsbasis, Lokalisierung und Kartierung, Navigation sowie Bedien- und Leitstandsebene belastbar arbeiten. Erst danach folgen Erweiterungen wie Vision oder Sprachschnittstelle.

\subsection{PF1 — Fahrkern und verteilte Echtzeitarchitektur}

Die erste Projektfrage untersucht, ob sich ein Fahrkern auf Basis einer verteilten ESP32-S3-Architektur so realisieren lässt, dass Motorregelung, Odometrie und die Verarbeitung von Geschwindigkeitskommandos deterministisch arbeiten. Die Validierung zeigt für die PID-Regelschleife einen Regeltakt von $50\,\mathrm{Hz}$ bei einem Jitter von unter $2\,\mathrm{ms}$. Der Datenverlust der micro-ROS-Kommunikation blieb unter $0{,}1\,\%$. Diese Messwerte belegen, dass die Trennung von Drive-Knoten und Sensor-Knoten blockierende Peripheriezugriffe vom Fahrkern fernhält und damit die zentrale Architekturabsicht erfüllt.

PF1 ist damit positiv beantwortet. Eine verteilte Echtzeitarchitektur mit UART-basierter micro-ROS-Anbindung kann den Fahrkern für einen kostengünstigen Innenraum-AMR deterministisch stabil betreiben, sofern die Aufgabenteilung zwischen Motorregelung und Sensorverarbeitung physisch getrennt bleibt.

\subsection{PF2 — Sensor- und Sicherheitsbasis sowie Lokalisierung und Kartierung}

Die zweite Projektfrage untersucht, welchen Beitrag Odometrie-Kalibrierung, IMU-Integration, Kanten-Erkennung und LiDAR-basierte Kartierung zur belastbaren Zustands- und Umgebungsbasis leisten. Die Validierung zeigt drei zentrale Ergebnisse. Erstens reduzierte die UMBmark-Kalibrierung den systematischen Odometriefehler um den Faktor $10$ bis $20$. Zweitens sank der laterale Drift auf einer Referenzfahrt über $1{,}0\,\mathrm{m}$ mit IMU-Korrektur von $3{,}6\,\mathrm{cm}$ auf $2{,}1\,\mathrm{cm}$, während der Gierfehler von $5{,}57^\circ$ auf $0{,}06^\circ$ sank. Drittens erreichte die Kartierung einen mittleren Absolute Trajectory Error von $0{,}16\,\mathrm{m}$ und blieb damit unter dem Akzeptanzkriterium von $0{,}20\,\mathrm{m}$. Zusätzlich reagierte die Sicherheitskette bei einem Kantenereignis in unter $50\,\mathrm{ms}$ (gemessen: $2{,}0\,\mathrm{ms}$).

PF2 ist damit weitgehend positiv beantwortet. Odometrie-Kalibrierung, IMU-Stützung, LiDAR-basierte Kartierung und priorisierte Sicherheitslogik bilden gemeinsam eine belastbare Grundlage für den Innenraumbetrieb. Die Einzelbewertung aller Sensorpfade einschließlich Ultraschall, Batterieüberwachung und IMU-Plausibilisierung ist mit $11/11$ bestandenen Testfällen abgeschlossen (vgl.\ Kapitel~6, Abschnitt~6.2.4).

\subsection{PF3 — Navigation sowie Bedien- und Leitstandsebene}

Die dritte Projektfrage untersucht, ob sich auf dieser Basis eine Systemarchitektur aufbauen lässt, die Zielanfahrten, Recovery-Verhalten, Telemetrie, manuelle Eingriffe und Audio-Rückmeldungen zusammenführt, ohne die Freigabelogik zu unterlaufen. Die Validierung zeigt für die Navigation eine mittlere Positionsabweichung von $6{,}4\,\mathrm{cm}$ und eine mittlere Winkelabweichung von $4{,}2^\circ$. Der Nav2-Stack umfuhr statische Hindernisse mit einem minimalen Passierabstand von $12\,\mathrm{cm}$. Die Recovery-Verfahren lösten $80\,\%$ der Blockaden eigenständig. In einer gezielt erzeugten Konfliktsituation überstimmte die Sicherheitslogik die gültige Bewegungsplanung, und das System kam $3\,\mathrm{cm}$ vor einer Absturzkante zum Stillstand. Diese Messwerte belegen, dass die Navigation funktional arbeitet und dass der sichere Zustand Vorrang vor dem zuletzt berechneten Bewegungsziel hat.

Für die Bedien- und Leitstandsebene liegt der Nachweis vor allem auf technischer Ebene vor. Telemetrie, Video und Diagnosefunktionen lassen sich betreiben, ohne den Fahrkern strukturell zu destabilisieren. Eine eigenständige Messreihe zur Bedienqualität, zur Verständlichkeit von Zustandsanzeigen oder zur Wirksamkeit von Audio-Rückmeldungen wurde jedoch nicht durchgeführt. PF3 ist deshalb im Kern positiv, in Bezug auf die Benutzerqualität jedoch nur teilweise beantwortet. Die Freigabelogik ist architektonisch richtig eingeordnet, aber noch nicht über alle Kommandoklassen systematisch vermessen.

\subsection{Kernaussage der Arbeit}

Die Arbeit zeigt, dass eine kostengünstige AMR-Architektur für strukturierte Innenräume technisch tragfähig ist, wenn der Fahrkern deterministisch ausgelegt, die Sensor- und Sicherheitsbasis priorisiert behandelt und die Bedien- und Leitstandsebene klar von der fahrkritischen Kette getrennt wird. Der belastbare Erkenntnisgewinn liegt damit weniger in einer einzelnen Demonstrationsfunktion als in der hierarchischen Kopplung der Systemebenen.

\section{Kritische Würdigung und Limitationen}

\subsection{Vollständigkeit des Nachweises}

Die Validierung bestätigt den fahrkritischen Kern, aber nicht alle Anforderungen mit gleicher Tiefe. Für den Fahrkern wurde die $360^\circ$-Rotation gemessen und ergab einen Ist-Wert von $358{,}1^\circ$ bei einem Restfehler von $1{,}88^\circ$ (vgl.\ Kapitel~6). Eine systematische Messreihe zum Nachlauf nach dem Stopp steht hingegen noch aus. In der Sensor- und Sicherheitsbasis sind alle $11/11$ Sensorik-Testfälle bestanden, darunter Kanten-Erkennung, Ultraschall, Batterieüberwachung und IMU-Plausibilisierung (vgl.\ Kapitel~6, Abschnitt~6.2.4). Für die Bedien- und Leitstandsebene fehlt eine eigenständige Untersuchung der Bedienqualität. Die Sprachschnittstelle ist in der Roadmap als Erweiterung beschrieben, gehört aber nicht zum validierten Kernumfang dieser Arbeit.

\subsection{Technische Grenzen des Prototyps}

Die Kalibrierung reduziert systematische Odometriefehler deutlich, kompensiert aber keinen Schlupf auf wechselnden oder unebenen Untergründen. Das Recovery-Verhalten erhöht die Robustheit der Navigation, beseitigt jedoch nicht alle Blockadesituationen. Die Vision-Pipeline erreicht mit dem Hailo-8L eine mittlere Inferenzzeit von rund $34\,\mathrm{ms}$ pro Bild und ist damit technisch leistungsfähig, bleibt jedoch außerhalb des sicherheitskritischen Kerns. Das Docking erreichte nach Optimierung der Dreifach-Bedingung eine Erfolgsquote von $100\,\%$ ($10/10$) und zeigte damit technische Machbarkeit als Erweiterungsfunktion. Beleuchtungsartefakte und ungünstige Ankunftsposen wirken sich weiterhin deutlich auf die monokulare Bildauswertung aus.

\subsection{Methodische Grenzen}

Das Vorgehen nach VDI 2206 strukturiert die Entwicklung des mechatronischen Gesamtsystems klar. Für die Softwareentwicklung mit wiederholten Anpassungen an ROS-2-Konfiguration, Recovery-Verhalten und Knotenstruktur entsteht jedoch ein iterativer Arbeitsmodus, der im linearen Bild des V-Modells nur eingeschränkt sichtbar wird. Die Arbeit bestätigt damit die Eignung des V-Modells als Ordnungsrahmen, nicht jedoch als vollständige Beschreibung realer Softwareiteration.

\subsection{Übertragbarkeit und wirtschaftliche Einordnung}

Die Ergebnisse gelten für strukturierte Innenräume mit kontrollierbaren Randbedingungen. Eine direkte Übertragung auf industrielle Daueranwendung, Mischverkehr mit Personen oder zertifizierungspflichtige Flurförderzeuge ist nicht zulässig. Dafür fehlen unter anderem redundante Sicherheitssensorik, formale Sicherheitsnachweise und eine robuste Hardwareintegration. Wirtschaftlich zeigt der Prototyp dennoch einen relevanten Punkt: Mit Consumer-Hardware und Open-Source-Software lässt sich ein funktional überzeugender Innenraum-AMR zu deutlich geringeren Kosten aufbauen als ein industrielles Serienprodukt. Die Differenz in Preis und Reifegrad markiert zugleich die Grenze zwischen Laborprototyp und industrieller Lösung.

\section{Ausblick und Weiterentwicklung}

\subsection{Kernpfad der Roadmap vervollständigen}

Der nächste Entwicklungsschritt sollte nicht mit zusätzlichen Demonstrationsfunktionen beginnen, sondern mit der vollständigen Schließung des Kernpfads. Dazu gehören eine erweiterte Messreihe für den Fahrkern, insbesondere zum Nachlaufverhalten nach dem Stopp, sowie eine formale Prüfung aller relevanten Zustandsübergänge der Freigabelogik. Für die Navigation ist insbesondere das Recovery-Verhalten weiter zu verbessern, bis Fehlfahrten, Sackgassen und lokale Blockaden reproduzierbar beherrscht werden. Diese Reihenfolge entspricht der Roadmap-Logik, in der quantitative Stabilität vor zusätzlicher Systemkomplexität steht.

\subsection{Bedien- und Leitstandsebene systematisch ausbauen}

Die Bedien- und Leitstandsebene ist technisch vorhanden, sollte aber als Betriebswerkzeug systematischer bewertet werden. Dafür sind Messgrößen für Telemetrie-Verzögerung, Zustandstransparenz, Eingriffszeit, Audio-Rückmeldungen und Fehlerdiagnose festzulegen. Ziel ist eine Benutzeroberfläche, die nicht nur funktioniert, sondern Betriebszustände eindeutig und belastbar vermittelt. Die Trennung zwischen Bedienung und Fahrlogik bleibt dabei unverändert, weil genau diese Trennung die Stabilität des Gesamtsystems absichert.

\subsection{Sprachschnittstelle als sichere Erweiterung integrieren}

Die Roadmap beschreibt die Sprachschnittstelle ausdrücklich als Erweiterung oberhalb des fahrkritischen Kerns. Diese Einordnung ist technisch zwingend. Sprachbefehle dürfen nicht direkt in rohe Geschwindigkeitsbefehle übergehen, sondern müssen über Intent-Erkennung, Freigabelogik und Missionskommando auf Navigation, Leitstand oder Audio wirken. Für die nächste Ausbaustufe bietet sich daher eine Architektur mit ReSpeaker Mic Array v2.0, \texttt{voice\_input\_node}, \texttt{speech\_to\_text\_node}, \texttt{voice\_intent\_node}, \texttt{voice\_command\_mux} und \texttt{text\_to\_speech\_node} an. Der erste validierbare Funktionsumfang sollte auf einen kleinen Wortschatz mit priorisiertem „Stopp", Moduswechseln und wenigen freigegebenen Missionskommandos begrenzt bleiben.

\subsection{Hardware- und Sicherheitsreife erhöhen}

Für einen robusteren Dauerbetrieb sind eine eigene Leiterplatine, definierte Steckverbinder, verbesserte elektromagnetische Verträglichkeit und eine mechanisch stabilere Sensorintegration naheliegende nächste Schritte. Zusätzlich sollte die Sicherheitskette um weitere Schutzpfade wie Bumper oder überwachte Not-Halt-Funktionen ergänzt werden. Für jeden Schritt in Richtung realer Betriebsumgebung steigt die Bedeutung normativer Anforderungen, insbesondere im Hinblick auf CE-Konformität und auf sicherheitsgerichtete Sensorik.

\section{Schluss}

Die Arbeit liefert den Nachweis, dass ein kostengünstiger Innenraum-AMR auf Basis einer verteilten Open-Source-Architektur technisch tragfähig aufgebaut und in seinen Kernfunktionen experimentell belegt werden kann. Der Fahrkern arbeitet deterministisch, die Sensor- und Sicherheitsbasis priorisiert sichere Zustände, Lokalisierung und Kartierung schaffen eine belastbare Umgebungsreferenz, und die Navigation erreicht eine praxistaugliche Zielanfahrt im Innenraum. Offene Punkte liegen vor allem in der vollständigen Abdeckung aller Akzeptanzkriterien und in der Robustheit erweiterter Interaktionsfunktionen. Genau daraus folgt der weitere Entwicklungsweg: zuerst den Kernpfad vollständig schließen, danach multimodale Erweiterungen wie Vision und Sprachschnittstelle kontrolliert und freigabelogisch abgesichert integrieren.
